\subsection{Arsitektur Sistem}

Arsitektur sistem audit trail yang diusulkan dapat dilihat pada Gambar~\ref{fig:arsitektur-sistem}. Sistem yang dirancang merupakan sebuah \textit{Audit Trail Service} yang diintegrasikan ke dalam arsitektur DecMed untuk menyediakan mekanisme pembangkitan, penyimpanan, dan pengambilan kembali (\textit{retrieval}) \textit{audit trail} sebagai bukti digital. Berbeda dengan komponen aplikasi pada DecMed yang berfokus pada penyediaan layanan Rekam Medis Elektronik (RME), sistem audit trail berfungsi sebagai infrastruktur pendukung yang mencatat aktivitas-aktivitas penting yang terjadi pada sistem sehingga dapat digunakan sebagai bukti dalam proses audit maupun investigasi.



Secara umum, arsitektur sistem terdiri atas empat bagian utama, yaitu komponen DecMed sebagai penghasil \textit{audit event}, \textit{Audit Trail Service}, \textit{Audit Repository}, dan auditor sebagai pihak yang melakukan pengambilan \textit{audit trail}. Komponen DecMed meliputi seluruh layanan dan aplikasi yang telah diidentifikasi pada tahap analisis sebagai sumber \textit{audit event}, seperti aplikasi klien, \textit{proxy re-encryption server}, \textit{smart contract}, maupun komponen lainnya. Setiap aktivitas yang memenuhi kriteria \textit{audit event} akan dikirimkan menuju \textit{Audit Trail Service} untuk diproses.

\textit{Audit Trail Service} merupakan komponen utama yang dikembangkan pada penelitian ini. Komponen ini bertanggung jawab untuk menerima \textit{audit event}, membangkitkan \textit{audit record} sesuai dengan skema yang telah ditentukan, serta menerapkan mekanisme yang diperlukan untuk menjaga integritas, keaslian (\textit{authenticity}), dan \textit{non-repudiation} dari setiap \textit{audit record}. Setelah diproses, \textit{audit record} akan diteruskan menuju repositori audit untuk disimpan sesuai dengan kebijakan retensi yang berlaku.

Penyimpanan \textit{audit trail} dilakukan pada \textit{Audit Repository} yang dirancang sebagai repositori khusus untuk menyimpan seluruh \textit{audit record}. Selain berfungsi sebagai media penyimpanan, repositori ini juga mendukung mekanisme pengelolaan retensi sehingga \textit{audit record} dapat dipertahankan sebagai bukti digital selama periode retensi yang telah ditentukan.

Untuk mendukung proses audit maupun investigasi, auditor dapat mengakses \textit{Audit Trail Service} melalui mekanisme \textit{query}. Berdasarkan parameter yang diberikan, sistem akan mengambil \textit{audit record} yang sesuai dari repositori audit dan mengembalikannya kepada auditor. Penelitian ini hanya berfokus pada mekanisme penyediaan kembali (\textit{retrieval}) \textit{audit trail} sebagai bukti digital. Fungsi lanjutan seperti pemantauan (\textit{monitoring}), analisis aktivitas, maupun pemberian peringatan (\textit{alerting}) tidak termasuk ke dalam ruang lingkup sistem yang dikembangkan.